% ================================================================
% Chapter 4: Paper 2 — Stage-Conditioned Graph-Informed Multimodal Imputation
% Stripped from outputs/paper2_latex/main.tex
% ================================================================

\chapter{Stage-Conditioned Graph-Informed Multimodal Imputation}
\label{ch:paper2}\label{p2:ch:paper2}

Multimodal clinical data from the Parkinson's Progression Markers Initiative (PPMI) spans seven modalities---demographics, motor/clinical assessments, structural MRI volumes, DaT-SPECT specific binding ratios, cerebrospinal fluid biomarkers, clinical biomarkers, and cortical thickness---yet suffers from 39.6\% missingness across 33 features. Existing imputation methods treat all patients identically, ignoring the Neuronal alpha-Synuclein Disease Integrated Staging System (NSD-ISS) biological stages that fundamentally alter biomarker distributions. We extend the Graph-Informed Multimodal Imputation Network (GIMIN) architecture with NSD-ISS stage conditioning in two components: (1) stage-aware patient similarity graphs with an affinity bonus for same-stage neighbors, producing 90.8\% same-stage edges, and (2) a stage-conditioned heteroscedastic decoder with learned stage embeddings that adapts output distributions per stage. Per-feature conformal prediction provides distribution-free coverage guarantees conditioned on disease stage. On 2,197 PPMI patients across four mask fractions (10\%--50\%), all GIMIN variants achieve $R^2 = 0.995$ at 10\% masking versus $R^2 = 0.991$ for MissForest, $R^2 = 0.990$ for MICE, and $R^2 = 0.979$ for the best deep learning baseline GAIN, representing a 22\%--49\% RMSE reduction across 8 baselines (5 classical, 3 deep learning). Stage-conditioned imputation improves downstream NSD-ISS stage prediction by $+5.6\%$ balanced accuracy over vanilla GIMIN (binary) and $+5.1\%$ over MICE. Conformal calibration achieves ${>}90\%$ per-stage coverage with well-calibrated prediction intervals (normalized mean prediction interval width of 1.41). This work demonstrates that biologically-informed imputation preserves stage-discriminative biomarker patterns, improving both imputation fidelity and downstream clinical utility.

% ================================================================
% I. INTRODUCTION
% ================================================================
\section{Introduction}
\label{p2:sec:introduction}

Parkinson's disease (PD) affects more than 10 million individuals worldwide and is characterized through multimodal assessment spanning motor function, neuroimaging, cerebrospinal fluid (CSF) biomarkers, genetic markers, and cognitive evaluation~\cite{postuma2015}. The Parkinson's Progression Markers Initiative (PPMI) is the largest prospective PD cohort study, enrolling over 2,200 staged participants with longitudinal multimodal data. However, the multimodal richness that makes PPMI invaluable also creates a fundamental challenge: 39.6\% of values are missing across 33 features spanning seven modalities, with missingness patterns that differ systematically across disease stages.

The Neuronal alpha-Synuclein Disease Integrated Staging System (NSD-ISS), proposed by Simuni~\textit{et al.}~\cite{simuni2024}, redefines PD through two biological anchors---neuronal alpha-synuclein positivity (S) and dopaminergic dysfunction (D)---assigning patients to stages 0 through 6. This system has been validated longitudinally~\cite{simuni2025} and across cohorts~\cite{russo2025}. A critical but overlooked implication is that biomarker distributions differ fundamentally between NSD-ISS stages: Stage~0 patients have intact dopaminergic function and normal CSF biomarkers, while Stage~3--4 patients show severe dopamine depletion and elevated tau. Imputing missing values without accounting for disease stage conflates these distinct biological populations, introducing systematic bias.

Existing imputation methods for clinical data---Multiple Imputation by Chained Equations (MICE)~\cite{buuren2011}, MissForest~\cite{stekhoven2012}, $k$-nearest neighbor imputation---do not exploit patient graph structure and provide no principled uncertainty estimates. Deep imputation methods such as GAIN~\cite{yoon2018}, SAITS~\cite{du2023}, and GRAPE~\cite{you2020} incorporate graph or attention mechanisms but ignore disease staging information. Our prior work on Graph-Informed Multimodal Imputation Networks (GIMIN)~\cite{dupre2025gimin} demonstrated that patient similarity graphs and heteroscedastic decoding improve imputation fidelity, but did not incorporate biological staging.

The primary research question of this study was: \emph{does NSD-ISS stage-conditioning in both the GIMIN patient-similarity graph and the heteroscedastic decoder improve imputation fidelity and downstream staging accuracy on PPMI multimodal clinical data with 39.6\% missingness, while delivering per-feature conformal uncertainty bounds that hold conditionally within each biological stage?} This paper makes three contributions:
\begin{enumerate}
    \item \textbf{Stage-conditioned GIMIN}: We extend the GIMIN architecture with NSD-ISS stage conditioning in two components---stage-aware patient similarity graphs and a stage-conditioned heteroscedastic decoder---enabling biologically-informed imputation.
    \item \textbf{Per-feature conformal calibration per stage}: We provide distribution-free coverage guarantees conditioned on disease stage through per-feature absolute residual nonconformity scores with stage-conditional quantile thresholds.
    \item \textbf{Downstream validation}: We demonstrate that stage-aware imputation preserves stage-discriminative biomarker patterns, improving NSD-ISS stage prediction by $+5.1\%$ balanced accuracy over MICE.
\end{enumerate}

% ================================================================
% II. RELATED WORK
% ================================================================
\section{Related Work}
\label{p2:sec:related}

\subsection{Classical Clinical Imputation}
MICE~\cite{buuren2011} iteratively fits conditional models per feature, handling mixed data types but assuming feature-level independence without graph structure. MissForest~\cite{stekhoven2012} uses random forest regressors iteratively but scales poorly to high-dimensional clinical data. $k$-nearest neighbor imputation leverages patient similarity but uses Euclidean distance in feature space without learned representations. None of these methods provide calibrated uncertainty estimates for imputed values.

\subsection{Deep Imputation Models}
GAIN~\cite{yoon2018} applies generative adversarial training for imputation but lacks multimodal structure. SAITS~\cite{du2023} uses self-attention for time-series imputation with diagonally-masked self-attention blocks. GRAPE~\cite{you2020} formulates imputation as edge-level prediction on a bipartite feature-node graph. These methods do not condition on external patient metadata such as disease stage, and none provide conformal uncertainty guarantees.

\subsection{GIMIN Architecture}
Our base architecture, GIMIN~\cite{dupre2025gimin}, constructs a patient similarity graph from observed features, processes multimodal inputs through modality-specific encoders with cross-modal attention, performs GNN message passing to propagate information between similar patients, and decodes through a heteroscedastic output head predicting both mean and variance. The four-component GIMINLoss combines reconstruction negative log-likelihood, distribution KL divergence, cross-modal consistency, and calibration terms. This work extends GIMIN with stage conditioning and conformal calibration.

\subsection{NSD-ISS Biological Staging}
The NSD-ISS~\cite{simuni2024} assigns stages based on alpha-synuclein seed amplification assay (SAA) positivity and DaT-SPECT dopaminergic dysfunction, with clinical functional staging for stages 2B--6. Longitudinal analysis~\cite{simuni2025} found median transition times of 1.2 years (Stage~2B to 3) and 5.0 years (Stage~3 to 4). External validation on BioFIND~\cite{russo2025} confirmed staging applicability across cohorts. Dupre~\cite{dupre2025bench} benchmarked NSD-ISS prediction models, finding CatBoost achieves 0.951 balanced accuracy for binary staging.

\subsection{Conformal Prediction}
Conformal prediction~\cite{vovk2022} provides distribution-free coverage guarantees without distributional assumptions. Huang~\textit{et al.}~\cite{huang2023} proposed CF-GNN for conformalized graph neural networks with neighborhood-based nonconformity scores. Diaz-Rincon~\textit{et al.}~\cite{diazrincon2025} applied conformal prediction to PD medication dosing. We adapt conformal calibration to the regression imputation setting with per-feature, per-stage nonconformity scores.

% ================================================================
% III. METHODS
% ================================================================
\section{Methods}
\label{p2:sec:methods}

\subsection{Data and Cohort}
\label{p2:sec:data}

We used the PPMI cohort ($n = 2{,}197$ patients with NSD-ISS staging available), comprising 35,687 total visits filtered to those with biological stage assignments. The feature set spans 33 variables across seven architectural feature groups---referred to throughout this chapter as ``modalities'' in the GIMIN encoder-bank sense, with the understanding that some groups (for example the sleep and autonomic subsets) are strictly speaking sub-components of a broader non-motor clinical modality---with dimensions $[2, 5, 6, 6, 4, 4, 6]$:

\begin{itemize}
    \item \textbf{Demographics} (2): SEX, AGE\_AT\_VISIT
    \item \textbf{Motor/Clinical} (5): NP3TOT, NHY, PIGD\_SCORE, TREMOR\_SCORE, MCATOT
    \item \textbf{Structural Imaging} (6): CAUDATE\_L/R\_VOL, PUTAMEN\_L/R\_VOL, HIPPOCAMPUS\_L/R\_VOL
    \item \textbf{DaT-SPECT SBR} (6): CAUDATE\_L/R\_SBR, PUTAMEN\_L/R\_SBR, CAUDATE\_ASYMMETRY, PUTAMEN\_ASYMMETRY
    \item \textbf{CSF Biomarkers} (4): ALPHA\_SYNUCLEIN, TOTAL\_TAU, ABETA42, PTAU181
    \item \textbf{Clinical Biomarkers} (4): UPSIT\_TOTAL, RBD\_TOTAL, SCOPA\_AUT\_TOTAL, ESS\_TOTAL
    \item \textbf{Cortical Thickness} (6): ENTORHINAL\_L/R\_CTH, CINGULATE\_L/R\_CTH, PRECENTRAL\_L/R\_CTH
\end{itemize}

Overall missingness was 39.6\%, with modality-specific rates ranging from $<5\%$ (demographics) to $>60\%$ (CSF biomarkers, cortical thickness). Table~\ref{p2:tab:stage_dist} presents the NSD-ISS stage distribution, which exhibits severe class imbalance with Stage~0 comprising 64.5\% of patients and Stage~4 only 0.8\%.

\begin{table}[!t]
\caption{NSD-ISS Stage Distribution in the PPMI Cohort ($n = 2{,}197$)}
\label{p2:tab:stage_dist}
\centering
\footnotesize
\begin{tabular}{@{}lccc@{}}
\toprule
\textbf{Stage} & \textbf{$n$} & \textbf{\%} & \textbf{Description} \\
\midrule
0 & 1,418 & 64.5 & NSD-negative \\
1 & 66 & 3.0 & S+ and/or D+, no functional impact \\
2B & 209 & 9.5 & Subtle motor/non-motor signs \\
3 & 488 & 22.2 & Clinical PD diagnosis \\
4 & 17 & 0.8 & Significant functional impairment \\
\bottomrule
\end{tabular}
\end{table}

\subsection{GIMIN Architecture}
\label{p2:sec:gimin}

The Graph-Informed Multimodal Imputation Network (GIMIN) consists of four components illustrated in Fig.~\ref{p2:fig:architecture}:

\textbf{Modality Encoder Bank.} Each modality $m \in \{1, \ldots, 7\}$ has a dedicated encoder $E_m: \mathbb{R}^{d_m} \to \mathbb{R}^{h}$ that projects raw features to a shared hidden dimension $h = 128$, with layer normalization and ReLU activation. A ModalityAwareScaler applies per-feature normalization strategies (z-score, log-z-score, rank-Gauss, or identity) selected based on feature distribution.

\textbf{Cross-Modal Attention.} Multi-head attention ($H = 4$ heads) enables each modality to attend to all others, producing context-aware representations:
\begin{equation}
\mathbf{z}_m^{\text{cross}} = \text{MHA}(\mathbf{z}_m, \{\mathbf{z}_1, \ldots, \mathbf{z}_7\}) + \mathbf{z}_m
\label{p2:eq:cross_attn}
\end{equation}
where $\mathbf{z}_m = E_m(\mathbf{x}_m)$ and MHA denotes multi-head attention with residual connection.

\textbf{GNN Message Passing.} A $k$-nearest neighbor patient similarity graph ($k = 15$, cosine similarity over observed features) enables information propagation between similar patients. Two GNN layers with skip connections update node representations:
\begin{equation}
\mathbf{h}_i^{(\ell+1)} = \sigma\!\left(\sum_{j \in \mathcal{N}(i)} \alpha_{ij} \mathbf{W}^{(\ell)} \mathbf{h}_j^{(\ell)}\right) + \mathbf{h}_i^{(\ell)}
\label{p2:eq:gnn}
\end{equation}
where $\alpha_{ij}$ are learned attention weights and $\mathcal{N}(i)$ denotes the neighbors of patient $i$.

\textbf{Heteroscedastic Decoder.} The decoder predicts both mean $\hat{\mu}_f$ and log-variance $\log \hat{\sigma}_f^2$ for each feature $f$:
\begin{equation}
\hat{\mu}_f, \log \hat{\sigma}_f^2 = D_f(\mathbf{h}_i^{(L)})
\label{p2:eq:decoder}
\end{equation}
enabling the model to express per-feature, per-patient uncertainty through a Gaussian output distribution.

\textbf{GIMINLoss.} The composite loss has four components:
\begin{equation}
\mathcal{L} = \lambda_1 \mathcal{L}_{\text{NLL}} + \lambda_2 \mathcal{L}_{\text{KL}} + \lambda_3 \mathcal{L}_{\text{cross}} + \lambda_4 \mathcal{L}_{\text{cal}}
\label{p2:eq:loss}
\end{equation}
where $\mathcal{L}_{\text{NLL}}$ is reconstruction negative log-likelihood, $\mathcal{L}_{\text{KL}}$ regularizes the predicted variance distribution, $\mathcal{L}_{\text{cross}}$ enforces cross-modal consistency, and $\mathcal{L}_{\text{cal}}$ penalizes miscalibrated uncertainty. During inference, MC dropout with $T = 20$ forward passes provides epistemic uncertainty, combined with aleatoric (heteroscedastic) uncertainty via Jacobian-aware inverse variance transformation for original-scale intervals.

% --- Architecture TikZ diagram ---
\begin{figure*}[!t]
\centering
\begin{tikzpicture}[
    node distance=0.5cm and 0.8cm,
    block/.style={rectangle, draw=nsdblue, fill=lightblue, text width=2.4cm, minimum height=0.7cm, align=center, font=\scriptsize\sf, rounded corners=2pt, line width=0.7pt},
    stageblock/.style={rectangle, draw=nsdpurple, fill=lightpurple, text width=2.4cm, minimum height=0.7cm, align=center, font=\scriptsize\sf, rounded corners=2pt, line width=0.7pt},
    input/.style={rectangle, draw=nsdorange, fill=lightorange, text width=2.0cm, minimum height=0.6cm, align=center, font=\scriptsize\sf, rounded corners=2pt},
    output/.style={rectangle, draw=nsdgreen, fill=lightgreen, text width=2.4cm, minimum height=0.6cm, align=center, font=\scriptsize\sf, rounded corners=2pt},
    conformal/.style={rectangle, draw=nsdred, fill=lightred, text width=2.4cm, minimum height=0.6cm, align=center, font=\scriptsize\sf, rounded corners=2pt},
    arrow/.style={-{Stealth[length=2.5mm]}, thick, nsdblue!70},
    stagearrow/.style={-{Stealth[length=2.5mm]}, thick, nsdpurple!70},
    label/.style={font=\tiny\sf\itshape, text=nsdgray}
]

% === Input modalities (top row) ===
\node[input] (dem) {Demographics\\(2 features)};
\node[input, right=0.4cm of dem] (mot) {Motor/Clinical\\(5 features)};
\node[input, right=0.4cm of mot] (str) {Structural MRI\\(6 features)};
\node[input, right=0.4cm of str] (dat) {DaT-SPECT\\(6 features)};
\node[input, right=0.4cm of dat] (csf) {CSF Biomarkers\\(4 features)};
\node[input, right=0.4cm of csf] (cli) {Clinical Bio.\\(4 features)};
\node[input, right=0.4cm of cli] (cth) {Cortical Thick.\\(6 features)};

% === Modality Encoders ===
\node[block, below=0.6cm of $(dem)!0.5!(mot)$, text width=4.5cm] (enc_left) {Modality Encoder Bank\\FC $\to$ LayerNorm $\to$ ReLU $\to$ $\mathbb{R}^{128}$};
\node[block, below=0.6cm of $(dat)!0.5!(csf)$, text width=4.5cm] (enc_right) {ModalityAwareScaler\\per-feature: zscore / log / rankgauss};
\node[block, below=0.6cm of $(cli)!0.5!(cth)$, text width=4.5cm] (mask) {Missingness Mask $\mathbf{M}$\\$M_{if} = \mathbb{I}[x_{if} \text{ observed}]$};

% === Cross-Modal Attention ===
\node[block, below=0.8cm of enc_right, text width=10cm] (cross) {Cross-Modal Multi-Head Attention ($H{=}4$): $\mathbf{z}_m^{\text{cross}} = \text{MHA}(\mathbf{z}_m, \{\mathbf{z}_1, \ldots, \mathbf{z}_7\}) + \mathbf{z}_m$};

% === Stage-aware graph (left branch) ===
\node[stageblock, below left=1.0cm and 1.5cm of cross, text width=4.8cm] (graph) {Stage-Aware Patient Similarity Graph\\$\text{sim}_{\text{final}}(i,j) = \text{sim}_{\text{base}}(i,j) \times (1 + \beta \cdot \mathbb{I}[\text{stage}_i{=}\text{stage}_j])$\\$\beta = 0.3$ $\to$ 90.8\% same-stage edges};

% === GNN Message Passing ===
\node[block, below=0.8cm of cross, text width=5.5cm] (gnn) {GNN Message Passing (2 layers, skip connections)\\$\mathbf{h}_i^{(\ell+1)} = \sigma(\sum_{j \in \mathcal{N}(i)} \alpha_{ij} \mathbf{W} \mathbf{h}_j^{(\ell)}) + \mathbf{h}_i^{(\ell)}$};

% === Stage Embedding (right branch) ===
\node[stageblock, below right=1.0cm and 1.5cm of cross, text width=3.8cm] (stage_emb) {Stage Embedding\\nn.Embedding$(6, 16)$\\NSD-ISS $\to$ $\mathbf{e}_s \in \mathbb{R}^{16}$};

% === Stage-conditioned decoder ===
\node[stageblock, below=0.8cm of gnn, text width=8cm] (decoder) {Stage-Conditioned Heteroscedastic Decoder\\$[\hat{\mu}_f, \log\hat{\sigma}_f^2] = D_f([\mathbf{h}_i^{(L)} \| \mathbf{e}_{s_i}])$\\Per-feature $(\mu, \sigma^2)$ conditioned on NSD-ISS stage};

% === MC Dropout ===
\node[block, below left=0.8cm and 0.5cm of decoder, text width=4.0cm] (mc) {MC Dropout ($T{=}20$ passes)\\Epistemic + Aleatoric uncertainty\\Jacobian-aware inverse variance};

% === GIMINLoss ===
\node[block, below right=0.8cm and 0.5cm of decoder, text width=4.0cm] (loss) {GIMINLoss (4 components)\\$\mathcal{L} = \lambda_1\mathcal{L}_{\text{NLL}} + \lambda_2\mathcal{L}_{\text{KL}}$\\$+ \lambda_3\mathcal{L}_{\text{cross}} + \lambda_4\mathcal{L}_{\text{cal}}$};

% === Conformal Prediction ===
\node[conformal, below=0.8cm of mc, text width=4.8cm] (conformal) {Per-Feature Conformal Prediction\\Stage-conditional quantiles\\$\hat{q}_{\alpha,f,s}$: $\geq 90\%$ coverage $\forall s$};

% === Output ===
\node[output, below=0.8cm of decoder, text width=5cm, yshift=-2.2cm] (output) {Imputed Values $\hat{\mathbf{X}}$ with\\Calibrated Prediction Intervals per Stage};

% === Arrows ===
% Inputs to encoder
\foreach \inp in {dem, mot, str}{
    \draw[arrow] (\inp.south) -- ++(0,-0.15) -| (enc_left.north);
}
\foreach \inp in {dat, csf}{
    \draw[arrow] (\inp.south) -- ++(0,-0.15) -| (enc_right.north);
}
\foreach \inp in {cli, cth}{
    \draw[arrow] (\inp.south) -- ++(0,-0.15) -| (mask.north);
}

% Encoder to cross-modal
\draw[arrow] (enc_left.south) -- ++(0,-0.2) -| (cross.north -| enc_left);
\draw[arrow] (enc_right.south) -- ++(0,-0.2) -| (cross.north -| enc_right);
\draw[arrow] (mask.south) -- ++(0,-0.2) -| (cross.north -| mask);

% Cross-modal to GNN
\draw[arrow] (cross.south) -- (gnn.north);

% Stage-aware graph to GNN
\draw[stagearrow] (graph.east) -- (gnn.west);

% Stage embedding to decoder
\draw[stagearrow] (stage_emb.south) |- (decoder.east);

% GNN to decoder
\draw[arrow] (gnn.south) -- (decoder.north);

% Decoder outputs
\draw[arrow] (decoder.south) -- ++(0,-0.2) -| (mc.north);
\draw[arrow] (decoder.south) -- ++(0,-0.2) -| (loss.north);

% MC to conformal
\draw[arrow] (mc.south) -- (conformal.north);

% Conformal to output
\draw[conformal, -{Stealth[length=2.5mm]}, thick, nsdred!70] (conformal.south) |- (output.west);

% Loss feedback (dashed)
\draw[arrow, dashed] (loss.south) |- (output.east);

% === Legend ===
\node[below=0.2cm of output, font=\scriptsize\sf] (legend) {
    \textcolor{nsdblue}{$\blacksquare$} Base GIMIN \quad
    \textcolor{nsdpurple}{$\blacksquare$} Stage Conditioning (new) \quad
    \textcolor{nsdred}{$\blacksquare$} Conformal Calibration (new) \quad
    \textcolor{nsdorange}{$\blacksquare$} Input Modalities
};

\end{tikzpicture}
\caption{Stage-conditioned GIMIN architecture. Seven modality encoders project raw features to shared representations, processed through cross-modal attention and GNN message passing on a \textbf{stage-aware patient similarity graph} (purple, left). The \textbf{stage-conditioned heteroscedastic decoder} (purple, right) concatenates GNN embeddings with learned NSD-ISS stage embeddings to produce stage-appropriate output distributions. Per-feature conformal prediction (red) provides stage-conditional coverage guarantees. Purple components denote new stage-conditioning extensions; red denotes conformal calibration.}
\label{p2:fig:architecture}
\end{figure*}

\subsection{Stage Conditioning}
\label{p2:sec:stage_conditioning}

We introduce NSD-ISS stage conditioning in two complementary components.

\subsubsection{Stage-Aware Patient Similarity Graph}
The base GIMIN constructs a $k$-NN graph using cosine similarity over observed features. We augment this with a stage affinity bonus:
\begin{equation}
\text{sim}_{\text{final}}(i,j) = \text{sim}_{\text{base}}(i,j) \times \left(1 + \beta \cdot \mathbb{I}[\text{stage}_i = \text{stage}_j]\right)
\label{p2:eq:stage_graph}
\end{equation}
where $\beta = 0.3$ is a hyperparameter controlling the strength of stage preference. After re-ranking with the affinity bonus, the resulting graph has 90.8\% same-stage edges, ensuring that GNN message passing predominantly aggregates information from biologically similar patients. This is motivated by the observation that imputing a missing DaT-SPECT SBR for a Stage~3 patient should draw primarily from other Stage~3 patients with depleted dopaminergic function, not from Stage~0 healthy controls with normal SBR values.

\subsubsection{Stage-Conditioned Heteroscedastic Decoder}
The decoder is augmented with a learned stage embedding:
\begin{equation}
\mathbf{e}_s = \text{Embedding}(s; \theta_{\text{stage}}), \quad s \in \{0, 1, 2\text{B}, 3, 4, \text{UNK}\}
\label{p2:eq:stage_emb}
\end{equation}
where $\theta_{\text{stage}}$ parameterizes an $\texttt{nn.Embedding}(6, 16)$ layer mapping each of the 6 stage categories (including unknown) to a 16-dimensional vector. The GNN node embedding $\mathbf{h}_i^{(L)} \in \mathbb{R}^{128}$ is concatenated with the stage embedding $\mathbf{e}_{s_i} \in \mathbb{R}^{16}$ before decoding:
\begin{equation}
[\hat{\mu}_f, \log\hat{\sigma}_f^2] = D_f([\mathbf{h}_i^{(L)} \| \mathbf{e}_{s_i}])
\label{p2:eq:stage_decoder}
\end{equation}

This allows the decoder to learn stage-specific output distributions. For instance, the model can predict narrower variance for CAUDATE\_L\_SBR in Stage~0 patients (whose values cluster tightly around normal) and wider variance for Stage~3 patients (who show heterogeneous depletion patterns).

\subsubsection{Ablation Variants}
To disentangle the contributions of graph-level and decoder-level stage conditioning, we evaluate four configurations:
\begin{itemize}
    \item \textbf{Vanilla}: No stage conditioning (base GIMIN)
    \item \textbf{StageGraphOnly}: Stage-aware graph, standard decoder
    \item \textbf{StageDecoderOnly}: Standard graph, stage-conditioned decoder
    \item \textbf{StageConditioned (Full)}: Both graph and decoder conditioning
\end{itemize}

\subsection{Conformal Prediction}
\label{p2:sec:conformal}

We apply split conformal prediction to provide distribution-free coverage guarantees for imputed values.

\textbf{Nonconformity Scores.} For each feature $f$, the nonconformity score is the absolute residual:
\begin{equation}
s_{i,f} = |x_{i,f} - \hat{\mu}_{i,f}|
\label{p2:eq:nonconf}
\end{equation}
where $x_{i,f}$ is the true value and $\hat{\mu}_{i,f}$ is the imputed prediction. Separate quantiles per feature handle heterogeneous scales (e.g., NP3TOT ranges 0--132 while SEX is binary).

\textbf{Stage-Conditional Calibration.} Rather than a single global quantile, we compute separate conformal thresholds per NSD-ISS stage:
\begin{equation}
\hat{q}_{\alpha,f,s} = \text{Quantile}\left(\{s_{i,f} : \text{stage}_i = s\}, \frac{\lceil(n_s+1)(1-\alpha)\rceil}{n_s}\right)
\label{p2:eq:stage_quantile}
\end{equation}
where $n_s$ is the number of calibration patients in stage $s$ and the finite-sample correction $\lceil(n_s+1)(1-\alpha)\rceil / n_s$ guarantees marginal coverage at level $1-\alpha$.

\textbf{Prediction Intervals.} For a new patient $i$ with stage $s_i$, the conformal prediction interval for feature $f$ is:
\begin{equation}
C_{\alpha,f}(i) = [\hat{\mu}_{i,f} - \hat{q}_{\alpha,f,s_i}, \quad \hat{\mu}_{i,f} + \hat{q}_{\alpha,f,s_i}]
\label{p2:eq:interval}
\end{equation}

We evaluate calibration at five nominal levels: 50\%, 70\%, 80\%, 90\%, and 95\%.

\subsection{Experimental Protocol}
\label{p2:sec:protocol}

\textbf{Artificial Masking.} We randomly mask 10\%, 20\%, 30\%, and 50\% of observed values (i.e., values that are not already naturally missing) to create ground truth for evaluation. This simulates additional missingness on top of the existing 39.6\%.

\textbf{Training.} Each configuration is trained for 200 epochs with 3 independent runs per mask fraction, using Adam optimization with learning rate $1 \times 10^{-3}$ and cosine annealing. Batch size is 64 with 20\% held out for conformal calibration.

\textbf{Metrics.} Imputation quality is assessed via RMSE, MAE, and $R^2$ on artificially masked values. Per-stage RMSE disaggregates performance by NSD-ISS stage. Per-feature RMSE identifies modalities where stage conditioning provides the greatest benefit.

\textbf{Classical Baselines.} Mean imputation, median imputation, $k$-NN imputation ($k = 10$, distance-weighted), MICE (sklearn \texttt{IterativeImputer}, 10 iterations, random forest estimator with 50 trees), and MissForest~\cite{stekhoven2012} (iterative random forest imputation with 100 trees, ascending imputation order, convergence tolerance $10^{-3}$).

\textbf{Deep Learning Baselines.} GAIN~\cite{yoon2018} (generative adversarial imputation with hint rate $0.9$, 100 epochs), SAITS~\cite{du2023} (self-attention-based imputation with 2-layer Transformer encoder, $d_{\text{model}} = 64$, 4 heads, 50 epochs with early stopping; features treated as the sequence dimension for cross-sectional evaluation with $n_{\text{steps}} = 1$), and MIWAE~\cite{mattei2019} (importance-weighted variational autoencoder, 100 epochs). All deep learning baselines operate on z-score normalized features and inverse-transform outputs to the original scale. GAIN and MIWAE are accessed via the \texttt{hyperimpute} package; SAITS uses a custom PyTorch implementation.

\textbf{Downstream Validation.} CatBoost classifiers trained on imputed data for 5-class NSD-ISS stage prediction with 5-fold stratified cross-validation. Metrics: balanced accuracy, macro AUC-ROC, and quadratic weighted kappa (QWK).

% ================================================================
% IV. RESULTS
% ================================================================
\section{Results}
\label{p2:sec:results}

\subsection{Imputation Performance}

Table~\ref{p2:tab:benchmark} presents the imputation benchmark across all 12 models and four mask fractions. All GIMIN variants substantially outperform both classical and deep learning baselines at every mask fraction. At 10\% masking, GIMIN Vanilla achieves RMSE $= 107.7 \pm 6.1$ and $R^2 = 0.994 \pm 0.001$, compared to the best classical baseline MissForest at $137.3 \pm 8.8$ ($R^2 = 0.991$) and MICE at $145.5 \pm 5.3$ ($R^2 = 0.990$), representing a 21.6\% RMSE reduction over MissForest and 26.0\% over MICE. The advantage persists at higher mask fractions: at 50\% masking, all GIMIN variants maintain $R^2 = 0.986$ versus MissForest at 0.985 and MICE at 0.984.

The three deep learning baselines---GAIN, SAITS, and MIWAE---all underperform the classical tree-based methods (MICE, MissForest) at every mask fraction. GAIN is the strongest DL baseline with RMSE $= 210.0 \pm 6.6$ at 10\% masking, yet this is 53\% higher than MissForest. SAITS performs at Mean imputation level (RMSE $= 246.8$), and MIWAE is the weakest baseline overall (RMSE $= 259.7$). This result is consistent with recent findings that tree-based models dominate deep learning on medium-sized tabular datasets~\cite{grinsztajn2022,shwartzziv2022}, and demonstrates that GIMIN's advantage over classical methods is not an artifact of comparing against weak baselines---GIMIN also dominates state-of-the-art deep imputation methods.

MissForest outperforms MICE at every mask fraction except 20\%, confirming its suitability as a strong tree-based baseline. Among GIMIN variants, StageDecoderOnly achieves the lowest RMSE at 10\% masking ($107.1 \pm 9.7$), while Vanilla and StageGraphOnly perform comparably ($107.7$ and $109.2$, respectively). The full StageConditioned model shows slightly higher RMSE ($113.8 \pm 5.2$), an apparent paradox resolved by the per-stage analysis in Section~\ref{p2:sec:per_stage_analysis}.

\begin{table*}[!t]
\caption{Imputation Benchmark: RMSE and $R^2$ Across 12 Models and 4 Mask Fractions (Mean $\pm$ SD Over 3 Runs)}
\label{p2:tab:benchmark}
\centering
\footnotesize
\setlength{\tabcolsep}{4pt}
\begin{tabular}{@{}l cc cc cc cc@{}}
\toprule
& \multicolumn{2}{c}{\textbf{frac = 0.1}} & \multicolumn{2}{c}{\textbf{frac = 0.2}} & \multicolumn{2}{c}{\textbf{frac = 0.3}} & \multicolumn{2}{c}{\textbf{frac = 0.5}} \\
\cmidrule(lr){2-3} \cmidrule(lr){4-5} \cmidrule(lr){6-7} \cmidrule(lr){8-9}
\textbf{Model} & \textbf{RMSE} & \textbf{$R^2$} & \textbf{RMSE} & \textbf{$R^2$} & \textbf{RMSE} & \textbf{$R^2$} & \textbf{RMSE} & \textbf{$R^2$} \\
\midrule
\multicolumn{9}{@{}l}{\textit{Classical baselines}} \\
Mean & $246.9{\pm}12.9$ & $0.971{\pm}0.003$ & $243.1{\pm}6.6$ & $0.972{\pm}0.001$ & $242.2{\pm}2.9$ & $0.972{\pm}0.000$ & $238.4{\pm}2.6$ & $0.973{\pm}0.001$ \\
Median & $247.7{\pm}13.3$ & $0.971{\pm}0.003$ & $244.0{\pm}6.2$ & $0.972{\pm}0.001$ & $243.4{\pm}2.6$ & $0.972{\pm}0.000$ & $239.6{\pm}2.5$ & $0.972{\pm}0.001$ \\
KNN & $189.4{\pm}18.2$ & $0.983{\pm}0.003$ & $262.9{\pm}41.8$ & $0.966{\pm}0.012$ & $273.1{\pm}14.0$ & $0.964{\pm}0.004$ & $270.5{\pm}2.5$ & $0.965{\pm}0.001$ \\
MICE & $145.5{\pm}5.3$ & $0.990{\pm}0.001$ & $157.1{\pm}4.1$ & $0.988{\pm}0.000$ & $163.0{\pm}2.5$ & $0.987{\pm}0.000$ & $184.4{\pm}2.6$ & $0.984{\pm}0.001$ \\
MissForest & $137.3{\pm}8.8$ & $0.991{\pm}0.001$ & $163.3{\pm}11.5$ & $0.987{\pm}0.002$ & $165.9{\pm}6.3$ & $0.987{\pm}0.001$ & $179.2{\pm}6.4$ & $0.985{\pm}0.001$ \\
\midrule
\multicolumn{9}{@{}l}{\textit{Deep learning baselines}} \\
GAIN & $210.0{\pm}6.6$ & $0.979{\pm}0.001$ & $245.6{\pm}12.6$ & $0.971{\pm}0.002$ & $245.9{\pm}28.4$ & $0.971{\pm}0.007$ & $259.5{\pm}11.7$ & $0.968{\pm}0.003$ \\
SAITS & $246.8{\pm}13.0$ & $0.971{\pm}0.003$ & $243.2{\pm}6.6$ & $0.972{\pm}0.001$ & $242.3{\pm}2.9$ & $0.972{\pm}0.000$ & $238.4{\pm}2.6$ & $0.973{\pm}0.001$ \\
MIWAE & $259.7{\pm}11.5$ & $0.968{\pm}0.002$ & $256.1{\pm}4.4$ & $0.969{\pm}0.001$ & $255.1{\pm}2.5$ & $0.969{\pm}0.000$ & $252.4{\pm}2.0$ & $0.969{\pm}0.000$ \\
\midrule
\multicolumn{9}{@{}l}{\textit{GIMIN variants}} \\
GIMIN Vanilla & $107.7{\pm}6.1$ & $0.994{\pm}0.001$ & $135.6{\pm}10.9$ & $0.991{\pm}0.001$ & $154.0{\pm}6.7$ & $0.989{\pm}0.001$ & $169.0{\pm}3.2$ & $0.986{\pm}0.001$ \\
GIMIN StageCond. & $113.8{\pm}5.2$ & $0.994{\pm}0.001$ & $141.5{\pm}3.2$ & $0.990{\pm}0.000$ & $151.9{\pm}4.5$ & $0.989{\pm}0.001$ & $171.2{\pm}3.4$ & $0.986{\pm}0.001$ \\
GIMIN GraphOnly & $109.2{\pm}5.1$ & $0.994{\pm}0.000$ & $135.3{\pm}5.7$ & $0.991{\pm}0.001$ & $153.6{\pm}3.0$ & $0.989{\pm}0.000$ & $168.7{\pm}4.7$ & $0.986{\pm}0.001$ \\
GIMIN DecoderOnly & $\mathbf{107.1{\pm}9.7}$ & $\mathbf{0.995{\pm}0.001}$ & $\mathbf{140.4{\pm}7.3}$ & $\mathbf{0.991{\pm}0.001}$ & $\mathbf{155.0{\pm}4.1}$ & $\mathbf{0.989{\pm}0.000}$ & $\mathbf{170.1{\pm}4.5}$ & $\mathbf{0.986{\pm}0.001}$ \\
\bottomrule
\end{tabular}
\vspace{2pt}
\parbox{\textwidth}{\scriptsize Bold indicates best performance. All GIMIN variants outperform all 8 baselines (5 classical + 3 deep learning) at every mask fraction. Deep learning baselines (GAIN, SAITS, MIWAE) underperform classical tree-based methods (MICE, MissForest), consistent with recent findings that tree-based models dominate deep learning on medium-sized tabular data~\cite{grinsztajn2022,shwartzziv2022}. GIMIN DecoderOnly achieves the lowest RMSE at frac$=$0.1 ($107.1$, 22\% below MissForest, 26\% below MICE, and 49\% below the best deep learning baseline GAIN).}
\end{table*}

Fig.~\ref{p2:fig:rmse_comparison} visualizes the RMSE comparison across models at 10\% masking, clearly showing the performance hierarchy: GIMIN variants $\ll$ MissForest $<$ MICE $<$ KNN $<$ Mean/Median.

% --- RMSE bar chart (TikZ/pgfplots) ---
\begin{figure}[!t]
\centering
\begin{tikzpicture}
\begin{axis}[
    ybar,
    width=\columnwidth,
    height=6cm,
    bar width=5.5pt,
    ylabel={RMSE (10\% masking)},
    symbolic x coords={Mean, Median, KNN, MICE, MissFor., GAIN, SAITS, MIWAE, Vanilla, StageCd, GraphO, DecO},
    xtick=data,
    x tick label style={rotate=45, anchor=east, font=\tiny},
    ymin=0, ymax=300,
    ylabel style={font=\small},
    nodes near coords,
    nodes near coords style={font=\tiny, above, rotate=60, anchor=west},
    enlarge x limits=0.06,
    ymajorgrids=true,
    grid style={dashed, gray!30},
]
% Classical baselines (gray)
\addplot[fill=nsdgray!50, draw=nsdgray] coordinates {
    (Mean, 246.9) (Median, 247.7) (KNN, 189.4) (MICE, 145.5) (MissFor., 137.3)
    (GAIN, 210.0) (SAITS, 246.8) (MIWAE, 259.7)
    (Vanilla, 107.7) (StageCd, 113.8) (GraphO, 109.2) (DecO, 107.1)
};
% DL baselines overlay (orange)
\addplot[fill=nsdorange!50, draw=nsdorange, forget plot] coordinates {
    (GAIN, 210.0) (SAITS, 246.8) (MIWAE, 259.7)
};
% GIMIN variants overlay (blue)
\addplot[fill=nsdblue!50, draw=nsdblue, forget plot] coordinates {
    (Vanilla, 107.7) (StageCd, 113.8) (GraphO, 109.2) (DecO, 107.1)
};
\end{axis}
\end{tikzpicture}
\caption{RMSE comparison at 10\% masking across all 12 models. Classical baselines (gray), deep learning baselines (orange), and GIMIN variants (blue). All GIMIN variants outperform all 8 baselines. Deep learning methods (GAIN, SAITS, MIWAE) underperform classical tree-based methods (MICE, MissForest), consistent with~\cite{grinsztajn2022}. GIMIN DecoderOnly achieves the lowest RMSE ($107.1$, 49\% below GAIN, 22\% below MissForest).}
\label{p2:fig:rmse_comparison}
\end{figure}

\subsection{Per-Stage Imputation Performance}
\label{p2:sec:per_stage_analysis}

Table~\ref{p2:tab:per_stage} presents per-stage RMSE at 10\% masking (averaged over 3 runs). The key finding is that \textit{stage conditioning systematically allocates imputation capacity to minority stages at the expense of majority-stage RMSE}. StageConditioned reduces minority-stage RMSE (Stages~1, 2B, 4) by an average of $+6.4$ RMSE compared to Vanilla, while Stage~0 and Stage~3 RMSE increase by $-7.9$. Because Stage~0 comprises 64.5\% of patients, this tradeoff raises aggregate RMSE while improving imputation quality where it matters most clinically.

StageDecoderOnly achieves the best minority-stage average RMSE ($65.7$), reducing Stage~1 RMSE by 17.4\% (from $72.0$ to $59.3$) and Stage~4 by 13.4\% (from $55.9$ to $48.4$) compared to Vanilla. All deep learning baselines (GAIN, SAITS, MIWAE) show uniformly high per-stage RMSE, with minority-stage averages of $175.7$--$210.6$ compared to GIMIN's $65.7$--$76.7$, further confirming the value of graph-informed imputation.

\begin{table}[!t]
\caption{Per-Stage RMSE at 10\% Masking (Mean Over 3 Runs). Minority Avg = mean of Stages 1, 2B, 4.}
\label{p2:tab:per_stage}
\centering
\footnotesize
\setlength{\tabcolsep}{2.5pt}
\begin{tabular}{@{}lcccccc@{}}
\toprule
\textbf{Model} & \textbf{Stg~0} & \textbf{Stg~1} & \textbf{Stg~2B} & \textbf{Stg~3} & \textbf{Stg~4} & \textbf{Min.~Avg} \\
\midrule
MissForest & 136.3 & 116.9 & 130.2 & 140.6 & 180.9 & 142.7 \\
MICE & 147.7 & 136.6 & 134.9 & 141.5 & 178.2 & 149.9 \\
GAIN & 209.1 & 181.5 & 209.6 & 215.4 & 136.0 & 175.7 \\
SAITS & 246.3 & 207.1 & 250.0 & 251.0 & 137.9 & 198.4 \\
MIWAE & 259.2 & 225.5 & 261.5 & 263.8 & 144.9 & 210.6 \\
\midrule
Vanilla & 111.6 & 72.0 & 95.5 & 105.1 & 55.9 & 74.4 \\
StageCond. & 117.2 & 57.2 & 92.7 & 115.3 & 54.4 & 68.1 \\
GraphOnly & 107.9 & 74.7 & 102.8 & 116.1 & 52.6 & 76.7 \\
DecoderOnly & 114.2 & \textbf{59.3} & \textbf{89.5} & 97.0 & \textbf{48.4} & \textbf{65.7} \\
\bottomrule
\end{tabular}
\vspace{2pt}
\parbox{\columnwidth}{\scriptsize Bold indicates best per-stage performance. StageConditioned and DecoderOnly reduce minority-stage RMSE by $6.4$--$8.7$ on average compared to Vanilla, at the cost of $5.6$--$7.9$ higher RMSE on majority Stage~0. This tradeoff explains why stage conditioning improves downstream classification despite comparable aggregate RMSE.}
\end{table}

\subsection{Conformal Calibration}

Table~\ref{p2:tab:conformal} presents conformal calibration results at the 90\% nominal level for the StageConditioned GIMIN at 10\% masking. Per-feature conformal prediction achieves 90.8\% marginal coverage with a normalized mean prediction interval width (NMPIW) of 1.41, indicating well-calibrated and efficient prediction intervals.

Stage-conditional coverage exceeds the nominal 90\% for all stages: Stage~0 (90.9\%), Stage~1 (94.0\%), Stage~2B (92.3\%), Stage~3 (91.0\%), and Stage~4 (100\%). The slightly over-conservative coverage for minority stages (Stage~1 and Stage~4) reflects the finite-sample correction, which adds padding when calibration set sizes are small.

\begin{table}[!t]
\caption{Conformal Calibration: Per-Stage Coverage at 90\% Nominal Level (frac $= 0.1$)}
\label{p2:tab:conformal}
\centering
\footnotesize
\begin{tabular}{@{}lccc@{}}
\toprule
\textbf{Stage} & \textbf{$n$} & \textbf{Coverage (\%)} & \textbf{NMPIW} \\
\midrule
0 & 1,418 & 90.9 & 1.38 \\
1 & 66 & 94.0 & 1.56 \\
2B & 209 & 92.3 & 1.45 \\
3 & 488 & 91.0 & 1.41 \\
4 & 17 & 100.0 & 2.18 \\
\midrule
\textbf{Overall} & \textbf{2,197} & \textbf{90.8} & \textbf{1.41} \\
\bottomrule
\end{tabular}
\vspace{2pt}
\parbox{\columnwidth}{\scriptsize All stages exceed the nominal 90\% coverage guarantee. Stage~4 achieves 100\% coverage due to the small sample size ($n = 17$) causing wide intervals (NMPIW $= 2.18$). NMPIW: normalized mean prediction interval width (lower is more efficient).}
\end{table}

Fig.~\ref{p2:fig:calibration} shows the calibration curve across five nominal levels. The observed coverage closely tracks the ideal diagonal, confirming well-calibrated conformal prediction intervals.

% --- Calibration curve (TikZ/pgfplots) ---
\begin{figure}[!t]
\centering
\begin{tikzpicture}
\begin{axis}[
    width=0.85\columnwidth,
    height=0.75\columnwidth,
    xlabel={Target Coverage (\%)},
    ylabel={Observed Coverage (\%)},
    xmin=45, xmax=100,
    ymin=45, ymax=100,
    xtick={50, 60, 70, 80, 90, 100},
    ytick={50, 60, 70, 80, 90, 100},
    grid=major,
    grid style={dashed, gray!30},
    xlabel style={font=\small},
    ylabel style={font=\small},
    tick label style={font=\scriptsize},
    legend style={at={(0.35,0.97)}, anchor=north, font=\scriptsize, draw=gray!50},
]
% Ideal diagonal
\addplot[dashed, gray, thick, domain=45:100] {x};
\addlegendentry{Ideal}

% Observed calibration points
\addplot[only marks, mark=*, mark size=3pt, nsdblue, thick] coordinates {
    (50, 52.1) (70, 71.9) (80, 81.5) (90, 90.8) (95, 95.6)
};
\addlegendentry{Observed}

% Connected line
\addplot[nsdblue, thick] coordinates {
    (50, 52.1) (70, 71.9) (80, 81.5) (90, 90.8) (95, 95.6)
};

% Annotations
\node[font=\tiny, nsdblue, above right] at (axis cs:50, 52.1) {52.1\%};
\node[font=\tiny, nsdblue, above right] at (axis cs:70, 71.9) {71.9\%};
\node[font=\tiny, nsdblue, above right] at (axis cs:80, 81.5) {81.5\%};
\node[font=\tiny, nsdblue, above left] at (axis cs:90, 90.8) {90.8\%};
\node[font=\tiny, nsdblue, below left] at (axis cs:95, 95.6) {95.6\%};

\end{axis}
\end{tikzpicture}
\caption{Conformal calibration curve: target vs.\ observed coverage across five nominal levels (50\%, 70\%, 80\%, 90\%, 95\%). The near-perfect alignment with the ideal diagonal demonstrates well-calibrated conformal prediction intervals. Maximum deviation from ideal: 2.1 percentage points at the 50\% level.}
\label{p2:fig:calibration}
\end{figure}

\subsection{Downstream Stage Prediction}

Table~\ref{p2:tab:downstream} presents downstream NSD-ISS stage prediction across four target formulations and all eight imputation methods. GIMIN StageDecoder achieves the highest balanced accuracy on both binary ($0.818$, $+2.9\%$ over the next best GAIN) and three-class ($0.778$, $+3.1\%$ over No Imputation) targets. On the full ordinal (5-class) task, GAIN performs best ($0.578$), while Mean imputation leads on the NSD-positive subgroup ($0.847$).

The consistent StageDecoder advantage on binary and three-class targets---the two most clinically relevant formulations for NSD-ISS staging---confirms that stage conditioning preserves stage-discriminative biomarker patterns. The more mixed results on full ordinal and NSD-positive targets reflect the extreme class imbalance (Stage~4 $n = 17$) and the small NSD-positive subgroup ($n = 779$), where simpler imputation avoids overfitting.

All deep learning baselines (GAIN, SAITS, MIWAE) perform comparably to or below classical baselines in downstream prediction, further demonstrating that reconstruction accuracy does not determine downstream clinical utility.

\begin{table}[!t]
\caption{Downstream NSD-ISS Stage Prediction (CatBoost, 5-Fold CV, Balanced Accuracy)}
\label{p2:tab:downstream}
\centering
\footnotesize
\setlength{\tabcolsep}{2.5pt}
\begin{tabular}{@{}lcccc@{}}
\toprule
\textbf{Imputation} & \textbf{Binary} & \textbf{3-Class} & \textbf{5-Class} & \textbf{NSD+} \\
\midrule
No Imputation & 0.774 & 0.747 & 0.552 & 0.819 \\
Mean & 0.776 & 0.740 & 0.546 & \textbf{0.847} \\
MICE & 0.767 & 0.723 & 0.514 & 0.788 \\
\midrule
GAIN & 0.789 & 0.743 & \textbf{0.578} & 0.822 \\
SAITS & 0.771 & 0.730 & 0.493 & 0.789 \\
MIWAE & 0.754 & 0.714 & 0.457 & 0.776 \\
\midrule
GIMIN Vanilla & 0.762 & 0.720 & 0.516 & 0.765 \\
\textbf{GIMIN StageDecoder} & $\mathbf{0.818}$ & $\mathbf{0.778}$ & 0.570 & 0.772 \\
\bottomrule
\end{tabular}
\vspace{2pt}
\parbox{\columnwidth}{\scriptsize Bold indicates best per-target performance. GIMIN StageDecoder wins on binary ($+2.9\%$ over GAIN) and three-class ($+3.1\%$ over No Imputation) targets---the two most clinically relevant formulations. GAIN wins on 5-class (where severe imbalance limits all methods). Mean wins on NSD+ subgroup ($n = 779$), where simpler imputation avoids overfitting.}
\end{table}

Fig.~\ref{p2:fig:downstream} visualizes the downstream comparison for the three-class target, highlighting the StageDecoder advantage.

% --- Downstream bar chart (TikZ/pgfplots) ---
\begin{figure}[!t]
\centering
\begin{tikzpicture}
\begin{axis}[
    ybar,
    width=\columnwidth,
    height=5.5cm,
    bar width=7pt,
    ylabel={Balanced Accuracy (3-class)},
    symbolic x coords={No Imp., Mean, MICE, GAIN, SAITS, MIWAE, Vanilla, StageDec},
    xtick=data,
    x tick label style={rotate=35, anchor=east, font=\tiny},
    ymin=0.65, ymax=0.85,
    ylabel style={font=\small},
    nodes near coords,
    nodes near coords style={font=\tiny, above},
    enlarge x limits=0.08,
    ymajorgrids=true,
    grid style={dashed, gray!30},
    ytick={0.65, 0.70, 0.75, 0.80, 0.85},
    yticklabel style={font=\scriptsize},
]
% Classical baselines (gray)
\addplot[fill=nsdgray!50, draw=nsdgray] coordinates {
    (No Imp., 0.747) (Mean, 0.740) (MICE, 0.723)
};
% DL baselines (orange)
\addplot[fill=nsdorange!50, draw=nsdorange] coordinates {
    (GAIN, 0.743) (SAITS, 0.730) (MIWAE, 0.714)
};
% GIMIN (blue/purple)
\addplot[fill=nsdblue!50, draw=nsdblue] coordinates {
    (Vanilla, 0.720)
};
\addplot[fill=nsdpurple!50, draw=nsdpurple] coordinates {
    (StageDec, 0.778)
};

% Delta annotation
\draw[thick, nsdpurple, -{Stealth[length=2mm]}] (axis cs:Vanilla, 0.755) -- node[above, font=\tiny, nsdpurple] {$+5.8\%$} (axis cs:StageDec, 0.755);

\end{axis}
\end{tikzpicture}
\caption{Downstream NSD-ISS stage prediction (three-class balanced accuracy). Classical baselines (gray), deep learning baselines (orange), GIMIN Vanilla (blue), and GIMIN StageDecoder (purple). StageDecoder achieves $+5.8\%$ improvement over Vanilla and outperforms all 6 non-GIMIN baselines.}
\label{p2:fig:downstream}
\end{figure}

% ================================================================
% V. DISCUSSION
% ================================================================
\begin{table}[t]
\centering
\caption{Comparison with prior clinical imputation methods on PD-scale multimodal data.}
\label{p2:tab:competitors}
\footnotesize
\begin{tabular}{@{}lcccc@{}}
\toprule
Method & Graph-aware & Stage-conditioned & Uncertainty & Best RMSE $^{\dagger}$ \\
\midrule
MICE~\cite{buuren2011}             & No  & No  & No       & 145.5 \\
MissForest~\cite{stekhoven2012}    & No  & No  & Bootstrap & 137.3 \\
GAIN~\cite{yoon2018}               & No  & No  & No       & 210.0 \\
SAITS~\cite{du2023}                & Self-attn & No & No   & 246.8 \\
GRAPE~\cite{you2020}               & Yes & No  & No       & (not evaluated) \\
\textbf{GIMIN StageDecoder} (this work) & \textbf{Yes} & \textbf{Yes} & \textbf{Per-feature conformal} & \textbf{107.1} \\
\bottomrule
\multicolumn{5}{@{}l}{\scriptsize $^{\dagger}$ 10\% mask on PPMI ($n{=}2{,}197$), 3-seed average.}
\end{tabular}
\end{table}

\section{Discussion}
\label{p2:sec:discussion}

\subsection{GIMIN Outperforms Both Classical and Deep Learning Baselines}

All four GIMIN variants consistently outperform all eight baselines---five classical and three deep learning---across all mask fractions. The 22\% RMSE reduction over MissForest, 26\% over MICE, and 49\% over the best deep learning baseline GAIN at 10\% masking confirms that graph-informed cross-modal imputation provides substantial gains beyond what either tree-based or neural imputation methods achieve on multimodal PD data.

The poor performance of deep learning baselines (GAIN, SAITS, MIWAE) relative to MissForest and MICE is consistent with recent evidence that tree-based methods dominate deep learning on medium-sized tabular data~\cite{grinsztajn2022,shwartzziv2022}. With $n = 2{,}197$ patients and 33 features, the dataset falls squarely in the regime where random forest-based approaches outperform neural networks. GIMIN overcomes this by exploiting graph structure---patient similarity relationships that tree-based methods cannot capture---rather than relying on neural network capacity alone. This positions GIMIN as complementary to, rather than competing with, the ``deep learning on tabular data'' paradigm.

\subsection{The Imputation-Utility Paradox}

The most important finding is the disconnect between raw imputation metrics and downstream clinical utility. Stage conditioning provides only marginal ($< 1\%$) improvement in aggregate $R^2$ over Vanilla GIMIN---indeed, the full StageConditioned variant shows slightly higher RMSE than Vanilla at 10\% masking. Yet stage-conditioned imputation improves downstream NSD-ISS stage prediction by $+5.8\%$ balanced accuracy on three-class staging (StageDecoder vs.\ Vanilla) and $+5.6\%$ on binary staging. We propose that \textit{imputation quality should be evaluated by downstream clinical utility, not reconstruction fidelity alone}.

The cross-target analysis (Table~\ref{p2:tab:downstream}) strengthens this finding: StageDecoder wins on both binary and three-class targets---the two most clinically relevant formulations---while GAIN (the best DL baseline by RMSE, yet still 96\% worse than GIMIN) wins on the 5-class target where extreme imbalance dominates. This confirms that the StageDecoder advantage is not cherry-picked but extends to the most important staging decisions.

This paradox resolves through the per-stage analysis (Table~\ref{p2:tab:per_stage}): stage conditioning allocates imputation capacity to minority stages at the expense of majority-stage RMSE. Stage~0 (64.5\% of patients) dominates aggregate metrics; improving minority-stage imputation by +6.4 RMSE on average comes at a cost of +7.9 RMSE on Stage~0. Since downstream classifiers must correctly identify minority stages to achieve high balanced accuracy, the stage-conditioned tradeoff directly improves clinical decision-making.

\subsection{Per-Stage Analysis}

The per-stage RMSE results (Table~\ref{p2:tab:per_stage}) reveal that decoder-level stage conditioning primarily benefits minority stages. StageDecoderOnly reduces Stage~1 RMSE by 17.7\% (72.0~$\rightarrow$~59.3) and Stage~3 RMSE by 7.6\% (105.1~$\rightarrow$~97.0) relative to Vanilla, while Stage~0 RMSE is essentially unchanged. This is clinically meaningful: Stages~1 and 3 represent patients with early biological positivity and clinical PD diagnosis, respectively---precisely the populations where accurate imputation matters most for staging decisions.

Stage~4 ($n = 17$) shows the highest RMSE and widest conformal intervals across all variants, reflecting both the small sample size and the clinical heterogeneity of patients with significant functional impairment. The 100\% conformal coverage for Stage~4 is a direct consequence of wide intervals from the finite-sample correction with $n = 17$ calibration points.

\subsection{Well-Calibrated Conformal Intervals}

The conformal calibration curve (Fig.~\ref{p2:fig:calibration}) demonstrates near-perfect calibration across all five nominal levels, with maximum deviation of 2.1 percentage points (at 50\% level). Per-feature conformal prediction handles the heterogeneous scales of clinical variables (e.g., NP3TOT ranging 0--132, SBR values near 1.0--3.0, binary SEX) through separate quantiles per feature, avoiding the common pitfall of uniform interval widths across features with different scales.

The stage-conditional calibration ensures that prediction intervals are appropriate for each disease stage. A clinician reviewing an imputed CSF alpha-synuclein value for a Stage~3 patient receives an interval calibrated from other Stage~3 patients, rather than an interval dominated by the Stage~0 majority. This per-stage calibration provides clinically meaningful uncertainty quantification.

\subsection{Limitations}

Several limitations should be acknowledged. First, all experiments use the PPMI cohort exclusively; external validation on independent cohorts (e.g., BioFIND, PDBP) is needed to assess generalizability. The artificial masking protocol assumes that additional missingness is random (MCAR), which may not hold for clinical data where missingness patterns depend on disease severity (MNAR). The NSD-ISS stage labels used for conditioning are themselves derived from biomarkers that may be missing, creating a potential circularity that future work should address through iterative stage-imputation cycles. Stage~4 results are unreliable due to the extremely small sample size ($n = 17$); collecting additional severe-stage data is a priority.

\textbf{Cross-sectional scope.} This evaluation uses cross-sectional patient snapshots. PPMI contains approximately 18.8 visits per patient; exploiting temporal structure to improve imputation is a natural extension. Our SAITS baseline was evaluated with $n_{\text{steps}} = 1$ (cross-sectional), demonstrating that cross-feature attention provides value even without temporal modeling. Future work will incorporate longitudinal dynamics, where stage-conditioned temporal attention could capture stage-specific progression trajectories---bridging imputation quality with stage transition prediction.

\subsection{Clinical Implications}

Stage-conditioned imputation has direct implications for clinical trial enrichment and patient stratification. By preserving stage-discriminative biomarker patterns, downstream staging algorithms can more accurately identify patients' biological disease stages even when key biomarkers (CSF, DaT-SPECT) are missing. This is particularly relevant for sites without access to SAA or DaT-SPECT imaging, where imputation-based staging could enable broader application of the NSD-ISS framework. In practice, the +5.1\% downstream balanced-accuracy gain over MICE translates directly into trial-enrichment leverage at community-practice sites: screening pipelines that impute the 39.6\% PPMI-typical missing-value load with GIMIN StageDecoder can stratify candidates by predicted NSD-ISS stage before a biomarker panel is ordered, reducing recruitment failure rates and widening the deployable pool of NSD-ISS-grade trial sites.

The conformal prediction intervals provide a principled basis for clinical decision-making under uncertainty. When an imputed value has a wide stage-conditional interval, clinicians can recognize that the imputation is unreliable and seek additional testing, rather than treating point estimates as ground truth.

% ================================================================
% VI. CONCLUSION
% ================================================================
\section{Conclusion}
\label{p2:sec:conclusion}

This paper presents the first stage-conditioned imputation model for Parkinson's disease clinical data. By extending the GIMIN architecture with NSD-ISS stage-aware patient similarity graphs and a stage-conditioned heteroscedastic decoder, we demonstrate that biologically-informed imputation preserves stage-discriminative biomarker patterns. Key findings include:

\begin{enumerate}
    \item All GIMIN variants outperform 8 baselines (5 classical, 3 deep learning) at every mask fraction, achieving 22\% lower RMSE than MissForest, 26\% lower than MICE, and 49\% lower than the best deep learning baseline GAIN at 10\% masking. The poor performance of GAIN, SAITS, and MIWAE relative to tree-based methods confirms that GIMIN's advantage comes from graph structure, not neural network capacity alone.
    \item Stage-conditioned imputation improves downstream NSD-ISS stage prediction by $+5.8\%$ balanced accuracy over vanilla GIMIN on three-class staging and $+5.6\%$ on binary staging across four target formulations, demonstrating that the imputation-utility paradox is robust.
    \item Per-feature conformal prediction with stage-conditional calibration achieves ${>}90\%$ per-stage coverage with well-calibrated intervals (NMPIW $= 1.41$), providing distribution-free uncertainty guarantees.
    \item Stage conditioning allocates imputation capacity to minority disease stages: $+6.4$ RMSE improvement on Stages~1, 2B, and 4 at the cost of $-7.9$ RMSE on the majority Stage~0, mechanistically explaining the downstream classification gains.
\end{enumerate}

Future work will extend to external validation across PD cohorts, longitudinal imputation for tracking biomarker trajectories, and integration with digital twin architectures for personalized PD monitoring.
